%%%%

%%%   Самарский государственный технический университет

%%%   Кафедра прикладной математики и информатики

%%%

%%%   Преамбула для статьи в журнал "Вестник Самарского государственного

%%%   технического университета. Серия: «Физико-математические науки»"

%%%   Ориентирована под MikTEx 2.4 for Win32

%%%

%%%   Хотя сам журнал собирается под GNU/Linux на дистрибутиве TeXLive



\documentclass[11pt,a4paper,twoside]{article}



\usepackage[cp1251]{inputenc}

\usepackage[T2A]{fontenc}

\usepackage[english,russian]{babel}

\usepackage{samgtubib}

\usepackage[dvips]{graphicx}

\usepackage[multiple]{footmisc}



\usepackage{samgtu}

\renewcommand{\VestnikYear}{2009} % Год издания Вестника

\renewcommand{\VestnikNumber}{2\,(19)} % Номер Вестника

\renewcommand{\VestnikMonth}{Ноябрь} % Месяц выхода Вестника

\renewcommand{\VestnikData}{01.11.09} % Дата подписания Вестника в печать

\renewcommand{\VestnikUPL}{26,64} % Усл. печ. л.

\renewcommand{\VestnikUIL}{25,73} % Уч.-изд. л.

\renewcommand{\VestnikRegNumber}{479/09} % Регистрационный номер

\renewcommand{\VestnikOrder}{994} % Номер заказа









%\newcommand{\mbf}{b}



%\begin{document}





\renewcommand{\firstpage}{157} %Начальная страница статьи для выноса в колонтитул

\renewcommand{\lastpage}{160} %Конечная страница статьи для выноса в колонтитул

%Расставляются при окончательной прогонке Вестника



%\Partition{Краткие сообщения}% Раздел Вестника к которому относится статья



%\ShortPartition{Механика деформируемого твёрдого тела}% Раздел Вестника к которому относится статья



\udk{539.376}





\title{Математическое моделирование длительной прочности толстостенных труб 

на основе концепции интегрально"=среднего эквивалентного 

напряжения}{Математическое моделирование длительной прочности толстостенных труб \dots} 



\engtitle{Mathematical modeling of long durability of thick--walled pipes on the basis 

of concept the integral--average equivalent pressure} 





\author{Е.\,В.~Башкинова}{Б\,а\,ш\,к\,и\,н\,о\,в\,а~Е.\,В.}

\engauthor{E.\,V.~Bashkinova}



\workabstract{Предложена обобщённая модель длительной прочности элементов конструкций на основе интегрально"=средних эквивалентных напряжений. Показано, что интегрально"=средние эквивалентные напряжения в процессе ползучести являются практически инвариантными величинами в процессе деформирования трубы от упругого (упругопластического) состояния вплоть до разрушения. Приведены результаты построения обобщённых моделей для толстостенных труб при ползучести для однопараметрического нагружения с использованием различных видов интегрально"=средних эквивалентных напряжений. Выполнена проверка адекватности данных расчёта экспериментальным данным.}



\engabstract{The generalized model of long durability construction's elements on the basis of integral--average equivalent stresses is offered. It is shown, that integral--average equivalent stresses are practically invariant sizes during deformation of a pipe from elastic (elastoplastic) conditions down to fracture. Results developed a generalized model for thick--walled pipes are resulted at creep for one--parametrical loading with use of various kinds integral--average equivalent stress are given. Check of adequacy calculation's data to experimental data is executed.

}



\date{12.12.2007} 



\address{\samgtu}

\engaddress{Samara State Technical University, Samara, Russia}



\email{}



\renewcommand{\baselinestretch}{.95}

\maketitle







Критерий длительной прочности на основе концепции эквивалентного 

напряжённого состояния позволяет установить эквивалентные напряжённые 

состояния, приводящие к разрушению за одно и тоже время $t_*$, а~также 

даёт возможность определить значение $t_*$ с~помощью самого простого 

испытания: при чистом растяжении. В работах [1,~2] было доказано, что 

величина интегрально"=средних эквивалентных напряжений в~процессе 

ползучести толстостенных труб (вплоть до разрушения) является близкой 

к~инвариантной при заданных внешних нагрузках и~при этом не зависит от 

показателя нелинейности установившейся ползучести $n$ в~достаточно широком 

диапазоне изменения величины $\beta = \frac ba$, где $b$ и~$a$ "--- внешний и~внутренний 

радиусы. Поэтому характеристики неоднородного напряжённого состояния, 

определяемые, например, в~результате решения задачи об установившейся 

ползучести толстостенной трубы, заменяются интегрально"=средними по 

поперечному сечению (радиусу) значениями. Таким образом, практически вводятся некоторые 

эффективные по объёму интегрально"=средние эквивалентные напряжения и~по 

отношению к~ним применяется концепция длительной прочности на основе 

эквивалентных напряжённых состояний. В~настоящей работе рассмотрены 

следующие виды эквивалентных напряжений: 

\begin{equation}

\label{bash:eq1}

\begin{array}{c}

\sigma _{\text  э\,1} = \sigma _{\max } ,\quad

\sigma _{\text  э\,2} = \sigma _{\max } - \sigma _{\min }, \\

\sigma _{\text  э\,3} = \frac{1}{2}\left( {\sigma _{\text  э\,1} + \sigma _{\text  э\,4} } \right),

\quad

\sigma _{\text  э\,4} = \frac{1}{\sqrt 2 }\sqrt {\left( {\sigma _r - \sigma _\theta } 

\right)^2 + \left( {\sigma _r - \sigma _z } \right)^2 + \left( {\sigma 

_\theta - \sigma _z } \right)^2} .

\end{array}

\end{equation}



Интегрально"=среднее значение эквивалентных напряжений рассчитывалось при 

осреднении по поперечному сечению или по радиусу по следующим формулам:



\vspace{-2mm}

\begin{equation}

\bar {\sigma }_{\text  э\,i}^1 = \frac{1}{b^2 - a^2}\int\limits_a^b {r\sigma _{\text  э\,i} \left( r \right)dr} ,

\quad

\bar {\sigma }_{\text э\,i}^2 = \frac{1}{b - a}\int\limits_a^b {\sigma _{\text э\,i} \left( 

r \right)dr}\quad  (i=1,\, 2,\, 3,\, 4).

\end{equation}

\vspace{-2mm}



Полученный анализ кинетики интегрально"=средних эквивалентных напряжений 

для элементов конструкций в условиях ползучести от начального упругого 

(упругопластического) состояния до состояния, соответствующего разрушению 

элемента конструкции, позволил сделать вывод о~неизменности величины 

среднего интегрального значения интенсивности напряжения для толстостенной 

трубы в процессе ползучести от момента $t=0$ вплоть до разрушения (с~погрешностью не более 1{,}5\,{\%}) для широкого диапазона изменения $n$ и~$\beta $.



Были проведены исследования задач для толстостенных трубы и~сферы под 

действием внутреннего давления, для бесконечно длинной трубы под действием 

внутреннего давления (гипотеза $\varepsilon _z = 0$ не выполняется), для 

толстостенной трубы при действии внутреннего давления и~растягивающей осевой 

нагрузки. Во всех случаях полученные результаты позволили сделать вывод, 

что, во"=первых, интегрально"=среднее значение эквивалентного напряжения 

можно рассчитывать по упругому (упругопластическому) решению, не привлекая 

для этого решение краевой задачи для установившейся ползучести, как это 

рекомендовалось ранее. Во"=вторых, длительную прочность толстостенной трубы 

можно спрогнозировать по любой одноосной теории ползучести и~длительной 

прочности, заменив в~соответствующей модели номинальное напряжение на 

интегрально"=среднее значение эквивалентного напряжения, рассчитанное по 

упругому (упругопластическому) решению при заданном режиме нагружения.



Если в~одноосном случае обозначить номинальное напряжение через 

$\sigma _0 $, интегрально"=среднее эквивалентное напряжение элемента конструкции при 

$t=0$ в~упругом (упругопластическом) состоянии через $\bar {\sigma }^y$, то 

возможны два варианта оценки длительной прочности. Первый связан с~расчётом 

$\bar {\sigma }^y$ и~проведением эксперимента над образцом при напряжении 

$\sigma _0 = \bar {\sigma }^y$. Второй вариант предполагает наличие 

одноосной теории ползучести и~длительной прочности. В~этом случае достаточно 

выполнить расчёт на основании любой одноосной модели при рассчитанном 

значении $\sigma _0 = \bar {\sigma }^y$. В~настоящей работе использовался 

энергетический вариант одноосной теории ползучести и~длительной прочности 

(без учёта первой стадии) вида [3]:

\begin{equation}

\label{bash:eq2}

\begin{array}{l}

\displaystyle 

\varepsilon = e + e^p + p,\qquad e = \frac\sigma E, 

\vspace{1mm}\\

\displaystyle 

\dot {e}^p = 

\left\{ {\begin{array}{l}

 0, \qquad \sigma \left( t \right) \leqslant \sigma _{\text{пр}}, \\ 

 \left\{ {\begin{array}{l}

 \lambda \left[ {a\left( {\sigma \left( t \right) - \sigma _{\text{пр}} } \right)^{n_1 } - e^p\left( t \right)} \right],\quad a\left( {\sigma \left( t \right) - \sigma _{\text{пр}} } \right)^{n_1 } > e^p\left( t \right), \\ 

 0, \qquad  a\left( {\sigma \left( t \right) - \sigma _{\text{пр}} } \right)^{n_1 } 

\leqslant e^p\left( t \right), \quad \sigma \left( t \right) > \sigma _{\text{пр}},

\end{array}} \right.  

 \end{array}} \right.

\vspace{1mm}\\

\displaystyle 

\dot {p} = c\left( {\frac{\sigma \left( t \right)}{\sigma _ * }} \right)^n,

\vspace{1mm}

\\

\displaystyle 

\sigma = \sigma _0 \left( {1 + \omega } \right),

\vspace{1mm}

\\

\displaystyle 

\dot {\omega } = \gamma \sigma \dot {e}^p + \alpha \sigma \dot {p},

\vspace{1mm}

\\

\displaystyle 

\gamma = \gamma _1 \left( {e^p} \right)^{m_2 },

\quad

\alpha = \alpha _1 \left( {\sigma _0 } \right)^{m_1 },

\end{array}

\end{equation}

где $\varepsilon $ "--- полная деформация; $e$ и~$e^p$ "--- упругая и~пластическая деформации, соответственно; 

$p$ "--- деформация ползучести; $E$ "--- модуль Юнга; $c$, $n$, $\sigma_* $, $\gamma_1 $, $m_2 $, $\alpha _1$, $m_1$, 

$a$, $\lambda$, $n_1 $ "--- константы модели; $\omega $ "--- параметр повреждённости;

$\sigma _0 $ и $\sigma $ "--- соответственно истинное и~номинальное напряжения. В~качестве критерия разрушения используется следующее интегральное соотношение энергетического типа:



\vspace{-2mm}

\begin{equation}

\label{bash:eq3}

\int\limits_0^{t_*} {\frac{\sigma \,de^p}{A_*^p }} + \int\limits_0^{t_*} 

\frac{\sigma \,dp}{A_*^c \left( {\sigma_0 } \right)} = 1,

\end{equation}

\vspace{-2mm}



\noindent

где $A_*^c = \alpha _A \left( {\sigma _0 } \right)^{m_A }$ ($\alpha _A $, 

$m_A $ и $A_ * ^p $ "--- постоянные характеристики материала), $t_*$ "---

время разрушения материала.



Модель (\ref{bash:eq2})--(\ref{bash:eq3}) при замене напряжения $\sigma _0 $ на интегрально"=среднее 

$\bar {\sigma }^y$ в~дальнейшем будем называть обобщённой моделью ползучести и~длительной прочности элемента конструкции. 



Были выполнены расчёты на основе обобщённой модели (\ref{bash:eq2})--(\ref{bash:eq3}) с~различными 

выражениями для эквивалентных напряжений, рассчитанных по формулам (\ref{bash:eq1}), (2), 

и~проведён анализ адекватности данных расчёта длительной прочности по 

обобщённой модели экспериментальным данным для толстостенных труб из 

стали 12ХМФ ($T=590\,{}^\circ \rm C$), стали 20 ($T=500\,{}^\circ \rm C$) и~сплава 

ЭИ694 ($T=700\,{}^\circ \rm C$), представленными в~[2, 4]. Кроме этого, был 

выполнен сравнительный анализ с~данными прямого численного решения краевой 

задачи ползучести и~разрушения толстостенной трубы [4].



Для сравнения результатов расчёта по модели (\ref{bash:eq2})--(\ref{bash:eq3}) с~экспериментальными 

данными и~данными численного решения краевой задачи для трубы использовалась 

величина среднего относительного отклонения данных расчёта от экспериментальных:



\vspace{-2mm}

$$

\Delta _i = \frac{1}{N}\sum\limits_{j = 1}^N {\left| {\frac{t_{ * \,j} - 

t_{ij} }{t_{ij} }} \right|}\quad  (i=1,\,2), 

$$

\vspace{-2mm}



\noindent

где $t_{ *\, j} $ "--- экспериментальные, $t_{ij} $ "--- расчётные данные ($j$ "--- номер реализации); 

$N$ "--- количество экспериментов ($i=1$ соответствует отклонению  расчётных данных, полученных по  обобщённой модели

(3)--(4); $i=2$ "--- отклонению  численного решения краевой задачи); 



\vspace{-2mm}

$$

\label{bash:eq4}

\Delta _3 = \frac{1}{N}\sum\limits_{j = 1}^N {\left| {\frac{t_{2j} - t_{1j} 

}{t_{1j} }} \right|} ,

$$

\vspace{-2mm}





\begin{wraptable}{O}{.675\textwidth}

\renewcommand{\bfdefault}{b}

\vspace{-1mm}

\centering

\caption{\small \bfseries Экспериментальные и расчётные значения времени до разрушения для труб из стали 12ХМФ 

($\bf \boldsymbol T=590\,{}^{\boldsymbol\circ} \rm\bf C$) при внутреннем давлении $\bf \boldsymbol q=28$~МПа}\vspace{-3mm}

\begin{tabular}

{c|c|c|c|c|c|c|c}

\hline

\No~п/п& 

$b$, мм& 

$a$, мм& 

$\beta$& 

$\bar {\sigma }_{\text{э}\,2}^y$, МПа& 

$t_ * $, ч& 

$t_1$, ч& 

$t_2$, ч \\

\hline

1& 

14,85& 

13,05& 

1,138& 

190,89& 

247& 

167,7& 

179,3 \\

%\hline

2& 

15,34& 

13,24& 

1,159& 

166,98& 

452& 

429,5& 

488,9 \\

%\hline

3& 

16,37& 

13,87& 

1,18& 

148,63& 

1131& 

915& 

1098,3 \\

%\hline

4& 

16,51& 

13,91& 

1,187& 

143,42& 

1463& 

1145& 

1371,1 \\

%\hline

5& 

16,68& 

13,98& 

1,22& 

123,32& 

1612& 

2895& 

1168,1 \\

\hline

\end{tabular}

\vspace{-3mm}

\renewcommand{\bfdefault}{bx}

\end{wraptable}



\noindent

где $t_{1j} $ и $t_{2j} $ "--- расчётные значения времени разрушения для 

того и~другого методов для $j$-ой реализации, $N$ "--- количество реализаций.



В качестве примера в табл.\,1 приведены экспериментальные и расчётные значения времени до 

разрушения при  $\bar {\sigma }_{\text э\, i}^2 $ для\,\, труб\, \, из\, стали\, \mbox{12ХМФ ($T=$} 



\begin{wraptable}{O}{.4\textwidth}

\renewcommand{\bfdefault}{b}

\vspace{-5mm}

\centering

\caption{\small \bfseries Значения среднего отклонения расчётных данных длительной прочности}\vspace{-3mm}

\begin{tabular}

{c|c|c|c}

\hline

Материал труб& 

$\Delta _1$,\,\%& 

$\Delta _2$,\,\%& 

$\Delta _3$,\,\%\\

\hline

Сталь 12ХМФ& 

29,64& 

18,60& 

24,04 \\

%\hline

Сталь 20& 

22,87& 

30,01& 

19,57 \\

%\hline

Сплав ЭИ694& 

37,89& 

41,56& 

6,67 \\

\hline

\end{tabular}





\vspace{3mm}

\centering

\caption{\small \bfseries Значения среднего отклонения расчётных данных длительной прочности

в~зависимости от выбора интегрально"=среднего значения  эквивалентных напряжений}\vspace{-3mm}

\begin{tabular}

{p{80pt}|p{12pt}|p{30pt}|p{30pt}}

\hline

~~~~Материал& 

$\sigma _{\text э\,i} $& 

$\Delta^*_1$,\,\% & 

$\Delta^*_2$,\,\% \\

\hline

\raisebox{-4.50ex}[0cm][0cm]{Сталь 12ХМФ}& 

$\sigma _{\text э\,1} $& 

20,38& 

16,77 \\

\cline{2-4} 

 & 

$\sigma _{\text э\,2} $& 

47,24& 

71,07 \\

\cline{2-4} 

 & 

$\sigma _{\text э\,3} $& 

20,87& 

21,79 \\

\cline{2-4} 

 & 

$\sigma _{\text э\,4} $& 

24,04& 

35,04 \\

\hline

\raisebox{-4.50ex}[0cm][0cm]{~~~~~Сталь 20}& 

$\sigma _{\text э\,1} $& 

22,87& 

30,01 \\

\cline{2-4} 

 & 

$\sigma _{\text э\,2} $& 

22,04& 

61,56 \\

\cline{2-4} 

 & 

$\sigma _{\text э\,3} $& 

49,66& 

47,27 \\

\cline{2-4} 

 & 

$\sigma _{\text э\,4} $& 

19,57& 

49,95 \\

\hline

\end{tabular}

\renewcommand{\bfdefault}{bx}

\vspace{-3mm}

\end{wraptable}



\noindent 

$=590\,{}^\circ \rm C$).

Здесь  $b$, $a$ "--- значения наружного и внутреннего радиусов, 

$\beta$ "--- их отношение;~  

$t_*$ "--- экспериментальные значения длительной прочности труб;

$t_1 $, $t_2 $ "--- расчётные значения, полученные  по обобщённой модели (\ref{bash:eq2})--(\ref{bash:eq3}) при $\bar {\sigma }_{\text э\,2}^y$ и на основе решения краевой задачи при внутреннем давлении $q=28$~МПа, соответственно.



В табл.\,2 приведены средние значения отклонения расчётных данных по 

обобщённой модели $\Delta _1$ и~численного решения краевой задачи $\Delta_2 $ от экспериментальных данных для толстостенных труб из~сталей 12ХМФ, стали~20 и~сплава ЭИ694; 

$\Delta _3 $ "--- среднее отклонение данных расчёта по обобщённой модели (\ref{bash:eq2})--(\ref{bash:eq3}) от данных численного решения краевой задачи. 

Видно, что величина $\Delta _3 $ имеет приемлемую (для процесса длительной прочности) величину того же порядка, что и~$\Delta _1 $ и~$\Delta _2 $. Отсюда следует, что величину интегрально"=среднего значения 

интенсивности напряжений $\bar {\sigma }_{\text э\,2}^2 $ можно рассчитывать по 

упругому (упругопластическому) решению и~использовать для оценки времени разрушения в~обобщённых теориях длительной прочности типа (\ref{bash:eq2})--(\ref{bash:eq3}), не прибегая к~решению краевых задач.







Таким образом, имея одноосную модель для материала и~упругое 

(упругопластическое) решение для напряжений можно достаточно надёжно 

прогнозировать длительную прочность элементов конструкций на основе 

интегрально"=средних напряжений, не прибегая к~решению краевых задач 

ползучести.







В табл.~3 приведены результаты подробного исследования соответствия данных 

расчёта длительной прочности по обобщённой модели (\ref{bash:eq2})--(\ref{bash:eq3}) при этом 

в~качестве интегрально"=средних напряжений принимались их  всевозможные реализации (\ref{bash:eq1}),~(2).

Здесь $\Delta^*_1$, $\Delta^*_2$  "--- величины средних отклонений расчётных данных длительной прочности от экспериментальных при $\bar {\sigma }_{\text  э\,i}^1$ и~$\bar {\sigma }_{\text  э\,i}^2$, соответственно.  



Приведенные результаты показывают, что осреднение по радиусу и поперечному сечению 

дают близкие (в~смысле математической статистики) погрешности. 

Однако в~определённом смысле осреднение по радиусу даёт всё же лучшие 

результаты, чем осреднение по~сечению.





\vspace{-3mm}





\begin{thebibliography}{4}



\item \textit{Башкинова, Е.\,В.} К обоснованию концепции <<эталонного>> напряжения в~задачах длительной 

прочности толстостенных трубы и сферы [Текст]~/  Е.\,В.~Башкинова~/ Актуальные проблемы современной 

науки: Тез. третьей международ. конф-ции молодых учёных. "--- Самара: СамГТУ, 2002. "--- Ч.~1. "--- С.~48--49.



\item \textit{Радченко, В.\,П.} Об одном подходе к оценке длительной прочности толстостенных труб на основе интегрально средних напряженных состояний 

[Текст]~/  В.\,П.~Радченко, Е.\,В.~Башкинова, С.\,Н.~Кубышкина~/\!/ Вестн. Сам. гос. техн. ун-та. Сер.: 

Физ.-мат. науки. "--- 2002. "--- \No~16. "--- С.~96--104. "--- ISBN~5--7964--0355--9.



\item \textit{Радченко, В.\,П.} Реологическое деформирование и~разрушение материалов и~элементов 

конструкций [Текст]~/ В.\,П.~Радченко, Ю.\,А.~Ерёмин. "--- М.: Машиностроение-1, 2004. "--- 265~с.

"--- ISBN~5--94275--111--0.



\item \textit{Радченко, В.\,П.} Математическая модель реологического деформирования и~разрушения 

толстостенной трубы [Текст]~/  В.\,П.~Радченко, С.\,Н.~Кубышкина~/\!/ 

Вестн. Сам. гос. техн. ун-та. Сер.: 

Физ.-мат. науки. "--- 1998. "--- \No~6. "--- С.~23--35. "--- ISBN~5--7964--0084--3.





\end{thebibliography}



\vspace{8mm}



\makeengtitle













%\end{document}

